% resume.tex -- LaTeX conversion of Michael_Lehn_Resume_SoftwareEngineer.pages
% Compile with: lualatex resume.tex   (or: make)
%
% Layout values are measured from the PDF exported by Pages, in PDF points
% (bp, 1/72 in; TeX's pt is slightly smaller):
%   US Letter, margins 45bp (left/right), text from 42.6bp (first baseline)
%   to 756bp; font Arial; body 10bp with Pages' single line spacing for
%   Arial (11.45bp).
%
% The public version omits the phone number. It is built from this file with
%   lualatex -jobname=resume_public '\def\PublicVersion{}% resume.tex -- LaTeX conversion of Michael_Lehn_Resume_SoftwareEngineer.pages
% Compile with: lualatex resume.tex   (or: make)
%
% Layout values are measured from the PDF exported by Pages, in PDF points
% (bp, 1/72 in; TeX's pt is slightly smaller):
%   US Letter, margins 45bp (left/right), text from 42.6bp (first baseline)
%   to 756bp; font Arial; body 10bp with Pages' single line spacing for
%   Arial (11.45bp).
%
% The public version omits the phone number. It is built from this file with
%   lualatex -jobname=resume_public '\def\PublicVersion{}% resume.tex -- LaTeX conversion of Michael_Lehn_Resume_SoftwareEngineer.pages
% Compile with: lualatex resume.tex   (or: make)
%
% Layout values are measured from the PDF exported by Pages, in PDF points
% (bp, 1/72 in; TeX's pt is slightly smaller):
%   US Letter, margins 45bp (left/right), text from 42.6bp (first baseline)
%   to 756bp; font Arial; body 10bp with Pages' single line spacing for
%   Arial (11.45bp).
%
% The public version omits the phone number. It is built from this file with
%   lualatex -jobname=resume_public '\def\PublicVersion{}% resume.tex -- LaTeX conversion of Michael_Lehn_Resume_SoftwareEngineer.pages
% Compile with: lualatex resume.tex   (or: make)
%
% Layout values are measured from the PDF exported by Pages, in PDF points
% (bp, 1/72 in; TeX's pt is slightly smaller):
%   US Letter, margins 45bp (left/right), text from 42.6bp (first baseline)
%   to 756bp; font Arial; body 10bp with Pages' single line spacing for
%   Arial (11.45bp).
%
% The public version omits the phone number. It is built from this file with
%   lualatex -jobname=resume_public '\def\PublicVersion{}\input{resume}'
\documentclass[10pt]{article}

\usepackage[letterpaper,top=32.6bp,bottom=36bp,left=45bp,right=45bp,
            headheight=0bp,headsep=0bp,footskip=0bp]{geometry}
\usepackage{fontspec}
\usepackage[HTML]{xcolor}
\usepackage{enumitem}
\usepackage{lua-ul}
\usepackage{hyperref}

% Pages breaks lines first-fit; TeX's paragraph optimizer would not
\directlua{dofile(kpse.find_file("firstfit.lua", "lua"))}

% Arial if installed, otherwise its metric-compatible clone Arimo.
% No ligatures, and no TeX input mapping (keeps straight quotes as typed).
\IfFontExistsTF{Arial}
  {\setmainfont{Arial}[Ligatures={CommonOff,TeXOff}]}
  {\setmainfont{Arimo}[Ligatures={CommonOff,TeXOff}]}

\definecolor{name}{HTML}{17365D}     % dark blue (name)
\definecolor{tagline}{HTML}{4F81BD}  % light blue (tagline, rule below name)
\definecolor{heading}{HTML}{365F91}  % muted blue (section headings)
\definecolor{rule}{HTML}{808080}     % gray section rules
\definecolor{link}{HTML}{0563C1}     % hyperlinks

\hypersetup{hidelinks, pdftitle={Michael Lehn -- Resume}, pdfauthor={Michael Lehn}}

% Hyperlink in the Pages style: blue and underlined (breakable across lines).
% The underline uses Arial's underline thickness (0.073em regular, 0.105em
% bold), centered 0.1em below the baseline.
\makeatletter
\def\link@thick@m{0.073em}
\def\link@thick@b{0.105em}
\def\link@thick@bx{0.105em}
\newunderlinetype\linkUnderLine[\number\dimexpr1em\relax\f@series]{%
  \leaders\vrule
    height \dimexpr-0.1em+0.5\dimexpr\csname link@thick@\f@series\endcsname\relax\relax
    depth  \dimexpr 0.1em+0.5\dimexpr\csname link@thick@\f@series\endcsname\relax\relax}
\makeatother
\newcommand{\link}[2]{\href{#1}{\textcolor{link}{\linkUnderLine{#2}}}}

% ---------------------------------------------------------------------------
% Font sizes (Pages single line spacing for Arial is 1.145 x font size;
% all body lines use the 10bp line pitch)
\newcommand{\namesize}{\fontsize{22bp}{25.2bp}\selectfont}
\newcommand{\taglinesize}{\fontsize{10.5bp}{16.8bp}\selectfont} % 16.8bp below the name
\newcommand{\headsize}{\fontsize{10.5bp}{11.45bp}\selectfont}
\newcommand{\bodysize}{\fontsize{10bp}{11.45bp}\selectfont}
\newcommand{\smallsize}{\fontsize{9.5bp}{11.45bp}\selectfont}
\newcommand{\contactsize}{\fontsize{9bp}{10.5bp}\selectfont}

% Separator in the header lines
\newcommand{\sep}{\ \ ·\ \ }

% ---------------------------------------------------------------------------
% Vertical spacing: every line uses the 11.45bp line pitch; the extra space
% between blocks is added with \addvspace, so adjacent amounts do not stack.
% The space above an entry depends on what precedes it (\cv@entryskip).
\makeatletter
\newcommand{\cv@entryskip}{4.6bp}

% Section heading: bold blue caps with a 0.5bp gray rule 2.9bp below the
% baseline
\newcommand{\cvsection}[1]{%
  \par\addvspace{6.85bp}%
  {\headsize\bfseries\color{heading}%
    \leavevmode\rlap{\color{rule}\rule[-3.15bp]{\linewidth}{0.5bp}}#1\par}%
  \nopagebreak\addvspace{2.2bp}%
  \gdef\cv@entryskip{4.8bp}%
}

% Entry: bold title, right-aligned small info (technologies or dates)
\newcommand{\cventry}[2]{%
  \par\addvspace{\cv@entryskip}%
  {\bfseries #1}\hfill{\smallsize #2}\kern0.85bp\par % Pages' right tab stop
  \nopagebreak
  \gdef\cv@entryskip{4.1bp}% next entry directly follows (education)
}

% Italic line below an entry (role, dissertation, ...)
\newcommand{\cvsubtitle}[1]{%
  \vspace{0.7bp}%
  {\smallsize\itshape #1\par}%
  \nopagebreak\addvspace{0.9bp}%
}

% Skills paragraph with a bold run-in label
\newcommand{\skill}[2]{%
  \textbf{#1} #2\par\addvspace{2bp}%
}

% Bullet list: bullet at 5bp, text at 15bp from the margin
\newenvironment{cvitems}
  {\begin{itemize}[label=•, leftmargin=15bp, labelwidth=10bp, labelsep=0bp,
                   align=left, topsep=0.55bp, partopsep=0bp, parsep=0bp,
                   itemsep=1.5bp, beginpenalty=10000]}
  {\end{itemize}\gdef\cv@entryskip{4.6bp}}
\makeatother

% ---------------------------------------------------------------------------
\setlength{\parindent}{0bp}
\setlength{\parskip}{0bp}
\setlength{\topskip}{10bp}
\pagestyle{empty}

\begin{document}
\bodysize
\hyphenpenalty=10000 % Pages does not hyphenate this document
\exhyphenpenalty=0
\tolerance=9999      % Pages stretches interword spaces freely in justified text
\emergencystretch=0bp
\frenchspacing       % uniform interword spaces after punctuation, as in Pages
\clubpenalty=10000   % no single lines of a paragraph or bullet at page breaks
\widowpenalty=10000
\makeatletter
\@clubpenalty=\clubpenalty % list items reset \clubpenalty to this value
\makeatother

% ===========================================================================
% Header
{\namesize\addfontfeatures{LetterSpace=1.2}\bfseries\color{name}%
  \rule{0bp}{22bp}% places the baseline 22bp below the top margin
  \rlap{\smash{\color{tagline}\rule[-6.5bp]{\linewidth}{1bp}}}%
  MICHAEL LEHN, Ph.D.\par}
{\centering
  {\taglinesize\addfontfeatures{LetterSpace=7.6}\color{tagline}%
    \spaceskip=\dimexpr\fontdimen2\font+0.8bp\relax % tracking applies to spaces, too
    Scientific Software Engineer\sep C/C++\sep Compilers \& LLVM\sep High-Performance Computing\par}
  \vspace{1.9bp}
  {\contactsize Ulm, Germany\sep Open to relocation to the United States\par}
  {\contactsize
    \link{mailto:michael.lehn@mac.com}{michael.lehn@mac.com}%
    \ifdefined\PublicVersion
      \ \ ·\ \link{https://www.linkedin.com/in/michael-l-2b4a4829b/}{LinkedIn}%
    \else
      \sep +1 (213) 232-7502\sep\link{https://www.linkedin.com/in/michael-l-2b4a4829b/}{LinkedIn}%
    \fi
    \sep\link{https://github.com/michael-lehn}{GitHub}%
    \ ·\ \ \link{https://www.uni-ulm.de/mawi/institut-fuer-numerische-mathematik/institut/mitarbeiter/mlehn/}{Ulm University}\par}
}
\vspace{-0.7bp}

% ===========================================================================
\cvsection{SUMMARY}

Software engineer and mathematician (Ph.D.) with 20+ years of hands-on experience developing performance-critical C/C++ software. Expertise spans numerical algorithms and high-performance linear algebra, memory hierarchy and cache-aware algorithms, SIMD, multithreaded and distributed computing, CPU architecture, compiler construction, and low-level systems programming. Recent projects include an LLVM-based systems programming language with native C interoperability, a minimal self-hosting compiler, and custom CPU architectures implemented on FPGA.\par

% ===========================================================================
\cvsection{TECHNICAL SKILLS}

\skill{Languages:}{C++ (modern, templates/compile-time techniques), C, Python, Assembly (x86-64, AArch64, custom RISC), LLVM IR, SystemVerilog, Java, Perl}
\skill{Performance \& HPC:}{Numerical linear algebra, BLAS, LAPACK, ScaLAPACK, SIMD micro-kernels (SSE3), cache-aware/blocked algorithms, std::thread, OpenMP, MPI, CUDA, GPU and distributed-memory computing, benchmarking and profiling}
\skill{Compilers \& Systems:}{LLVM, compiler construction, code generation, x86-64 System V ABI, C interoperability, ISA design, CPU architecture, FPGA (iCE40, ARTIQ/Sinara), bare-metal programming}
\skill{Tools:}{Linux/Unix, macOS, Git, Qt, CMake, GNU Make, GDB, Valgrind, perf, Intel VTune, NVIDIA Nsight}

% ===========================================================================
\cvsection{SELECTED PROJECTS (OPEN SOURCE)}

\cventry{{\mdseries\link{https://github.com/michael-lehn/ulmBLAS}{ulmBLAS}} — High-Performance BLAS Library}{C++ · Assembly}
\begin{cvitems}
  \item Designed and implemented a BLAS library based on the layered architecture pioneered by BLIS, combining portable blocking/packing algorithms in C++ with small, replaceable architecture-specific micro-kernels.
  \item Reached \textasciitilde97\% of Intel MKL DGEMM performance with BLIS micro-kernels; wrote custom SSE3 micro-kernels and a plug-in interface for externally optimized kernels.
  \item Drove blocked LAPACK routines (e.g., DGETRF) at performance comparable to vendor LAPACK implementations.
  \item Authored the \link{https://github.com/michael-lehn/gemm-tutorial}{GEMM Tutorial} on optimizing matrix multiplication from pure C to SSE micro-kernels; referenced by the \link{https://github.com/flame/how-to-optimize-gemm}{FLAME project's GEMM optimization resources}.
  \item Used in \link{https://www.meta.com/de/en/legal/quest/third-party-notices/14/}{Meta Quest software}; listed in Meta’s official third-party software notices.
\end{cvitems}

\cventry{{\mdseries\link{https://github.com/michael-lehn/abc-llvm}{ABC}} ("A Better C") — Systems Language \& LLVM Compiler}{C++ · LLVM}
\begin{cvitems}
  \item Designed a small systems language exposing C's machine model with simplified declaration syntax; built an LLVM-based compiler with interchangeable backends targeting LLVM IR and custom ULM assembly.
  \item Implemented native C interoperability with x86-64 System V ABI lowering, enabling direct use of existing C libraries without wrappers; demonstrated with raylib applications.
  \item Developed \link{https://github.com/michael-lehn/not-abc}{not-abc}, a minimal self-hosting compiler originally implemented in ABC, demonstrating the language's suitability for systems programming and compiler construction.
\end{cvitems}

\cventry{{\mdseries\link{https://github.com/michael-lehn/FLENS}{FLENS}} — Flexible Library for Efficient Numerical Solutions}{C++}
\begin{cvitems}
  \item Built a numerical computing framework combining MATLAB-like abstractions with direct access to optimized BLAS/LAPACK kernels, achieving zero abstraction overhead as verified by benchmarks and generated assembly.
  \item Designed a generic architecture allowing algorithms to operate unchanged on local, distributed, GPU-resident, and implicitly represented matrices.
  \item Used compile-time C++ techniques and expression analysis to map high-level operations directly to optimized numerical kernels and detect performance-critical temporaries.
\end{cvitems}

\cventry{{\mdseries\link{https://github.com/michael-lehn/ulm-on-ice}{ulm-on-iCE}} / {\mdseries\link{https://github.com/michael-lehn/ulm-generator}{ulm-generator}} — FPGA CPU \& ISA Toolchain}{SystemVerilog · C/C++ · Assembly}
\begin{cvitems}
  \item Designed and implemented a custom RISC CPU in SystemVerilog on an iCE40 FPGA, from digital logic and ISA design to assembly programs running on physical hardware.
  \item Defined a declarative ISA description format and built a generator producing an assembler, virtual machine, instruction-level debugger, and documentation from a single specification.
\end{cvitems}

% ===========================================================================
\cvsection{PROFESSIONAL EXPERIENCE}

\cventry{Ulm University — Institute of Numerical Mathematics · Ulm, Germany}{2008 – present}
\cvsubtitle{Tenured Faculty Member (Akademischer Rat auf Lebenszeit), Numerical Mathematics}
\begin{cvitems}
  \item Develop and maintain open-source numerical and systems software in C/C++: high-performance BLAS/LAPACK kernels (ulmBLAS), a numerical computing framework (FLENS), an LLVM-based systems language and compiler (ABC), and ISA/FPGA CPU tooling (ulm-generator / ulm-on-iCE).
  \item Own the full engineering lifecycle for these codebases: architecture, implementation, testing, performance benchmarking and profiling, documentation, and long-term maintenance.
  \item Design and teach courses in high-performance computing, compiler construction, systems programming, numerical mathematics, and scientific computing; the HPC curriculum progresses from memory hierarchy and cache-aware algorithms through SIMD micro-kernels and shared-memory parallelism to MPI and GPU computing.
  \item Built a full-stack computing curriculum spanning logic gates, registers and ALUs, ISA design, assembly, systems programming, and compiler construction, supported by custom-developed software and FPGA hardware.
\end{cvitems}

\cventry{Ulm Center for Scientific Computing (UZWR), Ulm University · Ulm, Germany}{2006 – 2008}
\cvsubtitle{Managing Director}
\begin{cvitems}
  \item Developed complete C++/Qt applications for industrial and research partners.
  \item Translated industrial and interdisciplinary partners' engineering and scientific problems into mathematical models and working numerical software; developed custom solutions where no suitable off-the-shelf software existed.
  \item Among numerous interdisciplinary projects, designed and developed a self-contained Windows application for MRI-based biomechanical analysis, implementing the complete software stack in-house—from FEM and numerical algorithms to data processing and the graphical user interface.
  \item Scoped and evaluated incoming projects, designed computational approaches, and either implemented solutions directly or coordinated them with researchers possessing the required expertise.
\end{cvitems}

\cventry{Ulm University — Institute of Numerical Mathematics · Ulm, Germany}{2002 – 2008}
\cvsubtitle{Research Associate}
\begin{cvitems}
  \item Conducted research and software development in scientific computing and numerical linear algebra, with a focus on high-performance numerical software in C++.
  \item Designed and developed FLENS (Flexible Library for Efficient Numerical Solutions) as the core of my doctoral research on efficient, abstraction-preserving numerical software in C++.
\end{cvitems}

\cventry{University of Southern California — Modeling and Simulation Laboratory · Los Angeles, CA}{1999 – 2001}
\cvsubtitle{Research Assistant / Teaching Assistant}
\begin{cvitems}
  \item Conducted research in applied mathematics and scientific computing with Prof. Chunming Wang in \link{https://dornsife.usc.edu/msl/faculty/}{USC's Modeling and Simulation Laboratory} while completing an M.S. in Applied Mathematics; master's research conducted in collaboration with NASA's Jet Propulsion Laboratory (JPL).
\end{cvitems}

% ===========================================================================
\cvsection{SELECTED INDUSTRY \& RESEARCH COLLABORATIONS}

\cventry{\link{https://www.dlr.de/de/qt}{German Aerospace Center (DLR)} — Real-Time Control / Quantum Technologies}{2025 – present}
\begin{cvitems}
  \item Collaborate with the DLR Institute of Quantum Technologies on control software for a space-ready quantum-memory system based on atomic-gas technology.
  \item Contribute to the architecture and development of the real-time control system and scientist-facing software, using ARTIQ/Sinara with an FPGA-based core device for deterministic control and hardware integration.
  \item Work on real-time state-machine control, experimental parameter management, measurement and imaging acquisition, hardware integration, and FPGA-controlled instrumentation.
\end{cvitems}

\cventry{\link{https://tensor-solutions.com}{Tensor AI Solutions} — HPC / Tensor Networks}{2024}
\begin{cvitems}
  \item Collaborated with Tensor AI Solutions on the design and development of Tartarus, a C++ HPC library for efficient tensor contractions in tree-tensor-network applications for machine learning and QML.
  \item Co-supervised a research-driven student project, contributing substantially to the numerical algorithms, performance-oriented design, and development of the library.
  \item Applied TTGT and Loops-over-GEMM (LoG) contraction schemes, BLAS kernels (GEMM, GEMV, GER, DOT), memory-layout-aware algorithms, and Python-based C++ code generation to optimize tensor contractions on CPUs.
\end{cvitems}

\cventry{\link{https://www.mercedes-benz-techinnovation.com/en}{Daimler TSS} (Daimler Group) — Machine Learning / Recurrent Neural Networks}{2017 – 2018}
\begin{cvitems}
  \item Collaborated closely with Daimler TSS on the mathematical analysis of recurrent neural networks, with a particular focus on training and the vanishing/exploding gradient problem.
  \item Co-supervised a research-driven bachelor's thesis, contributing substantially to the mathematical formulation, analysis, and development of the work.
  \item Applied regression and optimization methods, gradient-based training and backpropagation, and mathematical analysis of feed-forward and Elman recurrent neural networks.
\end{cvitems}

% ===========================================================================
\cvsection{EDUCATION}

\cventry{Ph.D., Numerical Mathematics — Ulm University, Germany}{2008}
\cvsubtitle{Dissertation: FLENS — A Flexible Library for Efficient Numerical Solutions · Advisor: Prof. Karsten Urban}

\cventry{M.S., Mathematics — Ulm University, Germany}{2002}

\cventry{M.S., Applied Mathematics — University of Southern California, Los Angeles, CA}{2001}

\end{document}
'
\documentclass[10pt]{article}

\usepackage[letterpaper,top=32.6bp,bottom=36bp,left=45bp,right=45bp,
            headheight=0bp,headsep=0bp,footskip=0bp]{geometry}
\usepackage{fontspec}
\usepackage[HTML]{xcolor}
\usepackage{enumitem}
\usepackage{lua-ul}
\usepackage{hyperref}

% Pages breaks lines first-fit; TeX's paragraph optimizer would not
\directlua{dofile(kpse.find_file("firstfit.lua", "lua"))}

% Arial if installed, otherwise its metric-compatible clone Arimo.
% No ligatures, and no TeX input mapping (keeps straight quotes as typed).
\IfFontExistsTF{Arial}
  {\setmainfont{Arial}[Ligatures={CommonOff,TeXOff}]}
  {\setmainfont{Arimo}[Ligatures={CommonOff,TeXOff}]}

\definecolor{name}{HTML}{17365D}     % dark blue (name)
\definecolor{tagline}{HTML}{4F81BD}  % light blue (tagline, rule below name)
\definecolor{heading}{HTML}{365F91}  % muted blue (section headings)
\definecolor{rule}{HTML}{808080}     % gray section rules
\definecolor{link}{HTML}{0563C1}     % hyperlinks

\hypersetup{hidelinks, pdftitle={Michael Lehn -- Resume}, pdfauthor={Michael Lehn}}

% Hyperlink in the Pages style: blue and underlined (breakable across lines).
% The underline uses Arial's underline thickness (0.073em regular, 0.105em
% bold), centered 0.1em below the baseline.
\makeatletter
\def\link@thick@m{0.073em}
\def\link@thick@b{0.105em}
\def\link@thick@bx{0.105em}
\newunderlinetype\linkUnderLine[\number\dimexpr1em\relax\f@series]{%
  \leaders\vrule
    height \dimexpr-0.1em+0.5\dimexpr\csname link@thick@\f@series\endcsname\relax\relax
    depth  \dimexpr 0.1em+0.5\dimexpr\csname link@thick@\f@series\endcsname\relax\relax}
\makeatother
\newcommand{\link}[2]{\href{#1}{\textcolor{link}{\linkUnderLine{#2}}}}

% ---------------------------------------------------------------------------
% Font sizes (Pages single line spacing for Arial is 1.145 x font size;
% all body lines use the 10bp line pitch)
\newcommand{\namesize}{\fontsize{22bp}{25.2bp}\selectfont}
\newcommand{\taglinesize}{\fontsize{10.5bp}{16.8bp}\selectfont} % 16.8bp below the name
\newcommand{\headsize}{\fontsize{10.5bp}{11.45bp}\selectfont}
\newcommand{\bodysize}{\fontsize{10bp}{11.45bp}\selectfont}
\newcommand{\smallsize}{\fontsize{9.5bp}{11.45bp}\selectfont}
\newcommand{\contactsize}{\fontsize{9bp}{10.5bp}\selectfont}

% Separator in the header lines
\newcommand{\sep}{\ \ ·\ \ }

% ---------------------------------------------------------------------------
% Vertical spacing: every line uses the 11.45bp line pitch; the extra space
% between blocks is added with \addvspace, so adjacent amounts do not stack.
% The space above an entry depends on what precedes it (\cv@entryskip).
\makeatletter
\newcommand{\cv@entryskip}{4.6bp}

% Section heading: bold blue caps with a 0.5bp gray rule 2.9bp below the
% baseline
\newcommand{\cvsection}[1]{%
  \par\addvspace{6.85bp}%
  {\headsize\bfseries\color{heading}%
    \leavevmode\rlap{\color{rule}\rule[-3.15bp]{\linewidth}{0.5bp}}#1\par}%
  \nopagebreak\addvspace{2.2bp}%
  \gdef\cv@entryskip{4.8bp}%
}

% Entry: bold title, right-aligned small info (technologies or dates)
\newcommand{\cventry}[2]{%
  \par\addvspace{\cv@entryskip}%
  {\bfseries #1}\hfill{\smallsize #2}\kern0.85bp\par % Pages' right tab stop
  \nopagebreak
  \gdef\cv@entryskip{4.1bp}% next entry directly follows (education)
}

% Italic line below an entry (role, dissertation, ...)
\newcommand{\cvsubtitle}[1]{%
  \vspace{0.7bp}%
  {\smallsize\itshape #1\par}%
  \nopagebreak\addvspace{0.9bp}%
}

% Skills paragraph with a bold run-in label
\newcommand{\skill}[2]{%
  \textbf{#1} #2\par\addvspace{2bp}%
}

% Bullet list: bullet at 5bp, text at 15bp from the margin
\newenvironment{cvitems}
  {\begin{itemize}[label=•, leftmargin=15bp, labelwidth=10bp, labelsep=0bp,
                   align=left, topsep=0.55bp, partopsep=0bp, parsep=0bp,
                   itemsep=1.5bp, beginpenalty=10000]}
  {\end{itemize}\gdef\cv@entryskip{4.6bp}}
\makeatother

% ---------------------------------------------------------------------------
\setlength{\parindent}{0bp}
\setlength{\parskip}{0bp}
\setlength{\topskip}{10bp}
\pagestyle{empty}

\begin{document}
\bodysize
\hyphenpenalty=10000 % Pages does not hyphenate this document
\exhyphenpenalty=0
\tolerance=9999      % Pages stretches interword spaces freely in justified text
\emergencystretch=0bp
\frenchspacing       % uniform interword spaces after punctuation, as in Pages
\clubpenalty=10000   % no single lines of a paragraph or bullet at page breaks
\widowpenalty=10000
\makeatletter
\@clubpenalty=\clubpenalty % list items reset \clubpenalty to this value
\makeatother

% ===========================================================================
% Header
{\namesize\addfontfeatures{LetterSpace=1.2}\bfseries\color{name}%
  \rule{0bp}{22bp}% places the baseline 22bp below the top margin
  \rlap{\smash{\color{tagline}\rule[-6.5bp]{\linewidth}{1bp}}}%
  MICHAEL LEHN, Ph.D.\par}
{\centering
  {\taglinesize\addfontfeatures{LetterSpace=7.6}\color{tagline}%
    \spaceskip=\dimexpr\fontdimen2\font+0.8bp\relax % tracking applies to spaces, too
    Scientific Software Engineer\sep C/C++\sep Compilers \& LLVM\sep High-Performance Computing\par}
  \vspace{1.9bp}
  {\contactsize Ulm, Germany\sep Open to relocation to the United States\par}
  {\contactsize
    \link{mailto:michael.lehn@mac.com}{michael.lehn@mac.com}%
    \ifdefined\PublicVersion
      \ \ ·\ \link{https://www.linkedin.com/in/michael-l-2b4a4829b/}{LinkedIn}%
    \else
      \sep +1 (213) 232-7502\sep\link{https://www.linkedin.com/in/michael-l-2b4a4829b/}{LinkedIn}%
    \fi
    \sep\link{https://github.com/michael-lehn}{GitHub}%
    \ ·\ \ \link{https://www.uni-ulm.de/mawi/institut-fuer-numerische-mathematik/institut/mitarbeiter/mlehn/}{Ulm University}\par}
}
\vspace{-0.7bp}

% ===========================================================================
\cvsection{SUMMARY}

Software engineer and mathematician (Ph.D.) with 20+ years of hands-on experience developing performance-critical C/C++ software. Expertise spans numerical algorithms and high-performance linear algebra, memory hierarchy and cache-aware algorithms, SIMD, multithreaded and distributed computing, CPU architecture, compiler construction, and low-level systems programming. Recent projects include an LLVM-based systems programming language with native C interoperability, a minimal self-hosting compiler, and custom CPU architectures implemented on FPGA.\par

% ===========================================================================
\cvsection{TECHNICAL SKILLS}

\skill{Languages:}{C++ (modern, templates/compile-time techniques), C, Python, Assembly (x86-64, AArch64, custom RISC), LLVM IR, SystemVerilog, Java, Perl}
\skill{Performance \& HPC:}{Numerical linear algebra, BLAS, LAPACK, ScaLAPACK, SIMD micro-kernels (SSE3), cache-aware/blocked algorithms, std::thread, OpenMP, MPI, CUDA, GPU and distributed-memory computing, benchmarking and profiling}
\skill{Compilers \& Systems:}{LLVM, compiler construction, code generation, x86-64 System V ABI, C interoperability, ISA design, CPU architecture, FPGA (iCE40, ARTIQ/Sinara), bare-metal programming}
\skill{Tools:}{Linux/Unix, macOS, Git, Qt, CMake, GNU Make, GDB, Valgrind, perf, Intel VTune, NVIDIA Nsight}

% ===========================================================================
\cvsection{SELECTED PROJECTS (OPEN SOURCE)}

\cventry{{\mdseries\link{https://github.com/michael-lehn/ulmBLAS}{ulmBLAS}} — High-Performance BLAS Library}{C++ · Assembly}
\begin{cvitems}
  \item Designed and implemented a BLAS library based on the layered architecture pioneered by BLIS, combining portable blocking/packing algorithms in C++ with small, replaceable architecture-specific micro-kernels.
  \item Reached \textasciitilde97\% of Intel MKL DGEMM performance with BLIS micro-kernels; wrote custom SSE3 micro-kernels and a plug-in interface for externally optimized kernels.
  \item Drove blocked LAPACK routines (e.g., DGETRF) at performance comparable to vendor LAPACK implementations.
  \item Authored the \link{https://github.com/michael-lehn/gemm-tutorial}{GEMM Tutorial} on optimizing matrix multiplication from pure C to SSE micro-kernels; referenced by the \link{https://github.com/flame/how-to-optimize-gemm}{FLAME project's GEMM optimization resources}.
  \item Used in \link{https://www.meta.com/de/en/legal/quest/third-party-notices/14/}{Meta Quest software}; listed in Meta’s official third-party software notices.
\end{cvitems}

\cventry{{\mdseries\link{https://github.com/michael-lehn/abc-llvm}{ABC}} ("A Better C") — Systems Language \& LLVM Compiler}{C++ · LLVM}
\begin{cvitems}
  \item Designed a small systems language exposing C's machine model with simplified declaration syntax; built an LLVM-based compiler with interchangeable backends targeting LLVM IR and custom ULM assembly.
  \item Implemented native C interoperability with x86-64 System V ABI lowering, enabling direct use of existing C libraries without wrappers; demonstrated with raylib applications.
  \item Developed \link{https://github.com/michael-lehn/not-abc}{not-abc}, a minimal self-hosting compiler originally implemented in ABC, demonstrating the language's suitability for systems programming and compiler construction.
\end{cvitems}

\cventry{{\mdseries\link{https://github.com/michael-lehn/FLENS}{FLENS}} — Flexible Library for Efficient Numerical Solutions}{C++}
\begin{cvitems}
  \item Built a numerical computing framework combining MATLAB-like abstractions with direct access to optimized BLAS/LAPACK kernels, achieving zero abstraction overhead as verified by benchmarks and generated assembly.
  \item Designed a generic architecture allowing algorithms to operate unchanged on local, distributed, GPU-resident, and implicitly represented matrices.
  \item Used compile-time C++ techniques and expression analysis to map high-level operations directly to optimized numerical kernels and detect performance-critical temporaries.
\end{cvitems}

\cventry{{\mdseries\link{https://github.com/michael-lehn/ulm-on-ice}{ulm-on-iCE}} / {\mdseries\link{https://github.com/michael-lehn/ulm-generator}{ulm-generator}} — FPGA CPU \& ISA Toolchain}{SystemVerilog · C/C++ · Assembly}
\begin{cvitems}
  \item Designed and implemented a custom RISC CPU in SystemVerilog on an iCE40 FPGA, from digital logic and ISA design to assembly programs running on physical hardware.
  \item Defined a declarative ISA description format and built a generator producing an assembler, virtual machine, instruction-level debugger, and documentation from a single specification.
\end{cvitems}

% ===========================================================================
\cvsection{PROFESSIONAL EXPERIENCE}

\cventry{Ulm University — Institute of Numerical Mathematics · Ulm, Germany}{2008 – present}
\cvsubtitle{Tenured Faculty Member (Akademischer Rat auf Lebenszeit), Numerical Mathematics}
\begin{cvitems}
  \item Develop and maintain open-source numerical and systems software in C/C++: high-performance BLAS/LAPACK kernels (ulmBLAS), a numerical computing framework (FLENS), an LLVM-based systems language and compiler (ABC), and ISA/FPGA CPU tooling (ulm-generator / ulm-on-iCE).
  \item Own the full engineering lifecycle for these codebases: architecture, implementation, testing, performance benchmarking and profiling, documentation, and long-term maintenance.
  \item Design and teach courses in high-performance computing, compiler construction, systems programming, numerical mathematics, and scientific computing; the HPC curriculum progresses from memory hierarchy and cache-aware algorithms through SIMD micro-kernels and shared-memory parallelism to MPI and GPU computing.
  \item Built a full-stack computing curriculum spanning logic gates, registers and ALUs, ISA design, assembly, systems programming, and compiler construction, supported by custom-developed software and FPGA hardware.
\end{cvitems}

\cventry{Ulm Center for Scientific Computing (UZWR), Ulm University · Ulm, Germany}{2006 – 2008}
\cvsubtitle{Managing Director}
\begin{cvitems}
  \item Developed complete C++/Qt applications for industrial and research partners.
  \item Translated industrial and interdisciplinary partners' engineering and scientific problems into mathematical models and working numerical software; developed custom solutions where no suitable off-the-shelf software existed.
  \item Among numerous interdisciplinary projects, designed and developed a self-contained Windows application for MRI-based biomechanical analysis, implementing the complete software stack in-house—from FEM and numerical algorithms to data processing and the graphical user interface.
  \item Scoped and evaluated incoming projects, designed computational approaches, and either implemented solutions directly or coordinated them with researchers possessing the required expertise.
\end{cvitems}

\cventry{Ulm University — Institute of Numerical Mathematics · Ulm, Germany}{2002 – 2008}
\cvsubtitle{Research Associate}
\begin{cvitems}
  \item Conducted research and software development in scientific computing and numerical linear algebra, with a focus on high-performance numerical software in C++.
  \item Designed and developed FLENS (Flexible Library for Efficient Numerical Solutions) as the core of my doctoral research on efficient, abstraction-preserving numerical software in C++.
\end{cvitems}

\cventry{University of Southern California — Modeling and Simulation Laboratory · Los Angeles, CA}{1999 – 2001}
\cvsubtitle{Research Assistant / Teaching Assistant}
\begin{cvitems}
  \item Conducted research in applied mathematics and scientific computing with Prof. Chunming Wang in \link{https://dornsife.usc.edu/msl/faculty/}{USC's Modeling and Simulation Laboratory} while completing an M.S. in Applied Mathematics; master's research conducted in collaboration with NASA's Jet Propulsion Laboratory (JPL).
\end{cvitems}

% ===========================================================================
\cvsection{SELECTED INDUSTRY \& RESEARCH COLLABORATIONS}

\cventry{\link{https://www.dlr.de/de/qt}{German Aerospace Center (DLR)} — Real-Time Control / Quantum Technologies}{2025 – present}
\begin{cvitems}
  \item Collaborate with the DLR Institute of Quantum Technologies on control software for a space-ready quantum-memory system based on atomic-gas technology.
  \item Contribute to the architecture and development of the real-time control system and scientist-facing software, using ARTIQ/Sinara with an FPGA-based core device for deterministic control and hardware integration.
  \item Work on real-time state-machine control, experimental parameter management, measurement and imaging acquisition, hardware integration, and FPGA-controlled instrumentation.
\end{cvitems}

\cventry{\link{https://tensor-solutions.com}{Tensor AI Solutions} — HPC / Tensor Networks}{2024}
\begin{cvitems}
  \item Collaborated with Tensor AI Solutions on the design and development of Tartarus, a C++ HPC library for efficient tensor contractions in tree-tensor-network applications for machine learning and QML.
  \item Co-supervised a research-driven student project, contributing substantially to the numerical algorithms, performance-oriented design, and development of the library.
  \item Applied TTGT and Loops-over-GEMM (LoG) contraction schemes, BLAS kernels (GEMM, GEMV, GER, DOT), memory-layout-aware algorithms, and Python-based C++ code generation to optimize tensor contractions on CPUs.
\end{cvitems}

\cventry{\link{https://www.mercedes-benz-techinnovation.com/en}{Daimler TSS} (Daimler Group) — Machine Learning / Recurrent Neural Networks}{2017 – 2018}
\begin{cvitems}
  \item Collaborated closely with Daimler TSS on the mathematical analysis of recurrent neural networks, with a particular focus on training and the vanishing/exploding gradient problem.
  \item Co-supervised a research-driven bachelor's thesis, contributing substantially to the mathematical formulation, analysis, and development of the work.
  \item Applied regression and optimization methods, gradient-based training and backpropagation, and mathematical analysis of feed-forward and Elman recurrent neural networks.
\end{cvitems}

% ===========================================================================
\cvsection{EDUCATION}

\cventry{Ph.D., Numerical Mathematics — Ulm University, Germany}{2008}
\cvsubtitle{Dissertation: FLENS — A Flexible Library for Efficient Numerical Solutions · Advisor: Prof. Karsten Urban}

\cventry{M.S., Mathematics — Ulm University, Germany}{2002}

\cventry{M.S., Applied Mathematics — University of Southern California, Los Angeles, CA}{2001}

\end{document}
'
\documentclass[10pt]{article}

\usepackage[letterpaper,top=32.6bp,bottom=36bp,left=45bp,right=45bp,
            headheight=0bp,headsep=0bp,footskip=0bp]{geometry}
\usepackage{fontspec}
\usepackage[HTML]{xcolor}
\usepackage{enumitem}
\usepackage{lua-ul}
\usepackage{hyperref}

% Pages breaks lines first-fit; TeX's paragraph optimizer would not
\directlua{dofile(kpse.find_file("firstfit.lua", "lua"))}

% Arial if installed, otherwise its metric-compatible clone Arimo.
% No ligatures, and no TeX input mapping (keeps straight quotes as typed).
\IfFontExistsTF{Arial}
  {\setmainfont{Arial}[Ligatures={CommonOff,TeXOff}]}
  {\setmainfont{Arimo}[Ligatures={CommonOff,TeXOff}]}

\definecolor{name}{HTML}{17365D}     % dark blue (name)
\definecolor{tagline}{HTML}{4F81BD}  % light blue (tagline, rule below name)
\definecolor{heading}{HTML}{365F91}  % muted blue (section headings)
\definecolor{rule}{HTML}{808080}     % gray section rules
\definecolor{link}{HTML}{0563C1}     % hyperlinks

\hypersetup{hidelinks, pdftitle={Michael Lehn -- Resume}, pdfauthor={Michael Lehn}}

% Hyperlink in the Pages style: blue and underlined (breakable across lines).
% The underline uses Arial's underline thickness (0.073em regular, 0.105em
% bold), centered 0.1em below the baseline.
\makeatletter
\def\link@thick@m{0.073em}
\def\link@thick@b{0.105em}
\def\link@thick@bx{0.105em}
\newunderlinetype\linkUnderLine[\number\dimexpr1em\relax\f@series]{%
  \leaders\vrule
    height \dimexpr-0.1em+0.5\dimexpr\csname link@thick@\f@series\endcsname\relax\relax
    depth  \dimexpr 0.1em+0.5\dimexpr\csname link@thick@\f@series\endcsname\relax\relax}
\makeatother
\newcommand{\link}[2]{\href{#1}{\textcolor{link}{\linkUnderLine{#2}}}}

% ---------------------------------------------------------------------------
% Font sizes (Pages single line spacing for Arial is 1.145 x font size;
% all body lines use the 10bp line pitch)
\newcommand{\namesize}{\fontsize{22bp}{25.2bp}\selectfont}
\newcommand{\taglinesize}{\fontsize{10.5bp}{16.8bp}\selectfont} % 16.8bp below the name
\newcommand{\headsize}{\fontsize{10.5bp}{11.45bp}\selectfont}
\newcommand{\bodysize}{\fontsize{10bp}{11.45bp}\selectfont}
\newcommand{\smallsize}{\fontsize{9.5bp}{11.45bp}\selectfont}
\newcommand{\contactsize}{\fontsize{9bp}{10.5bp}\selectfont}

% Separator in the header lines
\newcommand{\sep}{\ \ ·\ \ }

% ---------------------------------------------------------------------------
% Vertical spacing: every line uses the 11.45bp line pitch; the extra space
% between blocks is added with \addvspace, so adjacent amounts do not stack.
% The space above an entry depends on what precedes it (\cv@entryskip).
\makeatletter
\newcommand{\cv@entryskip}{4.6bp}

% Section heading: bold blue caps with a 0.5bp gray rule 2.9bp below the
% baseline
\newcommand{\cvsection}[1]{%
  \par\addvspace{6.85bp}%
  {\headsize\bfseries\color{heading}%
    \leavevmode\rlap{\color{rule}\rule[-3.15bp]{\linewidth}{0.5bp}}#1\par}%
  \nopagebreak\addvspace{2.2bp}%
  \gdef\cv@entryskip{4.8bp}%
}

% Entry: bold title, right-aligned small info (technologies or dates)
\newcommand{\cventry}[2]{%
  \par\addvspace{\cv@entryskip}%
  {\bfseries #1}\hfill{\smallsize #2}\kern0.85bp\par % Pages' right tab stop
  \nopagebreak
  \gdef\cv@entryskip{4.1bp}% next entry directly follows (education)
}

% Italic line below an entry (role, dissertation, ...)
\newcommand{\cvsubtitle}[1]{%
  \vspace{0.7bp}%
  {\smallsize\itshape #1\par}%
  \nopagebreak\addvspace{0.9bp}%
}

% Skills paragraph with a bold run-in label
\newcommand{\skill}[2]{%
  \textbf{#1} #2\par\addvspace{2bp}%
}

% Bullet list: bullet at 5bp, text at 15bp from the margin
\newenvironment{cvitems}
  {\begin{itemize}[label=•, leftmargin=15bp, labelwidth=10bp, labelsep=0bp,
                   align=left, topsep=0.55bp, partopsep=0bp, parsep=0bp,
                   itemsep=1.5bp, beginpenalty=10000]}
  {\end{itemize}\gdef\cv@entryskip{4.6bp}}
\makeatother

% ---------------------------------------------------------------------------
\setlength{\parindent}{0bp}
\setlength{\parskip}{0bp}
\setlength{\topskip}{10bp}
\pagestyle{empty}

\begin{document}
\bodysize
\hyphenpenalty=10000 % Pages does not hyphenate this document
\exhyphenpenalty=0
\tolerance=9999      % Pages stretches interword spaces freely in justified text
\emergencystretch=0bp
\frenchspacing       % uniform interword spaces after punctuation, as in Pages
\clubpenalty=10000   % no single lines of a paragraph or bullet at page breaks
\widowpenalty=10000
\makeatletter
\@clubpenalty=\clubpenalty % list items reset \clubpenalty to this value
\makeatother

% ===========================================================================
% Header
{\namesize\addfontfeatures{LetterSpace=1.2}\bfseries\color{name}%
  \rule{0bp}{22bp}% places the baseline 22bp below the top margin
  \rlap{\smash{\color{tagline}\rule[-6.5bp]{\linewidth}{1bp}}}%
  MICHAEL LEHN, Ph.D.\par}
{\centering
  {\taglinesize\addfontfeatures{LetterSpace=7.6}\color{tagline}%
    \spaceskip=\dimexpr\fontdimen2\font+0.8bp\relax % tracking applies to spaces, too
    Scientific Software Engineer\sep C/C++\sep Compilers \& LLVM\sep High-Performance Computing\par}
  \vspace{1.9bp}
  {\contactsize Ulm, Germany\sep Open to relocation to the United States\par}
  {\contactsize
    \link{mailto:michael.lehn@mac.com}{michael.lehn@mac.com}%
    \ifdefined\PublicVersion
      \ \ ·\ \link{https://www.linkedin.com/in/michael-l-2b4a4829b/}{LinkedIn}%
    \else
      \sep +1 (213) 232-7502\sep\link{https://www.linkedin.com/in/michael-l-2b4a4829b/}{LinkedIn}%
    \fi
    \sep\link{https://github.com/michael-lehn}{GitHub}%
    \ ·\ \ \link{https://www.uni-ulm.de/mawi/institut-fuer-numerische-mathematik/institut/mitarbeiter/mlehn/}{Ulm University}\par}
}
\vspace{-0.7bp}

% ===========================================================================
\cvsection{SUMMARY}

Software engineer and mathematician (Ph.D.) with 20+ years of hands-on experience developing performance-critical C/C++ software. Expertise spans numerical algorithms and high-performance linear algebra, memory hierarchy and cache-aware algorithms, SIMD, multithreaded and distributed computing, CPU architecture, compiler construction, and low-level systems programming. Recent projects include an LLVM-based systems programming language with native C interoperability, a minimal self-hosting compiler, and custom CPU architectures implemented on FPGA.\par

% ===========================================================================
\cvsection{TECHNICAL SKILLS}

\skill{Languages:}{C++ (modern, templates/compile-time techniques), C, Python, Assembly (x86-64, AArch64, custom RISC), LLVM IR, SystemVerilog, Java, Perl}
\skill{Performance \& HPC:}{Numerical linear algebra, BLAS, LAPACK, ScaLAPACK, SIMD micro-kernels (SSE3), cache-aware/blocked algorithms, std::thread, OpenMP, MPI, CUDA, GPU and distributed-memory computing, benchmarking and profiling}
\skill{Compilers \& Systems:}{LLVM, compiler construction, code generation, x86-64 System V ABI, C interoperability, ISA design, CPU architecture, FPGA (iCE40, ARTIQ/Sinara), bare-metal programming}
\skill{Tools:}{Linux/Unix, macOS, Git, Qt, CMake, GNU Make, GDB, Valgrind, perf, Intel VTune, NVIDIA Nsight}

% ===========================================================================
\cvsection{SELECTED PROJECTS (OPEN SOURCE)}

\cventry{{\mdseries\link{https://github.com/michael-lehn/ulmBLAS}{ulmBLAS}} — High-Performance BLAS Library}{C++ · Assembly}
\begin{cvitems}
  \item Designed and implemented a BLAS library based on the layered architecture pioneered by BLIS, combining portable blocking/packing algorithms in C++ with small, replaceable architecture-specific micro-kernels.
  \item Reached \textasciitilde97\% of Intel MKL DGEMM performance with BLIS micro-kernels; wrote custom SSE3 micro-kernels and a plug-in interface for externally optimized kernels.
  \item Drove blocked LAPACK routines (e.g., DGETRF) at performance comparable to vendor LAPACK implementations.
  \item Authored the \link{https://github.com/michael-lehn/gemm-tutorial}{GEMM Tutorial} on optimizing matrix multiplication from pure C to SSE micro-kernels; referenced by the \link{https://github.com/flame/how-to-optimize-gemm}{FLAME project's GEMM optimization resources}.
  \item Used in \link{https://www.meta.com/de/en/legal/quest/third-party-notices/14/}{Meta Quest software}; listed in Meta’s official third-party software notices.
\end{cvitems}

\cventry{{\mdseries\link{https://github.com/michael-lehn/abc-llvm}{ABC}} ("A Better C") — Systems Language \& LLVM Compiler}{C++ · LLVM}
\begin{cvitems}
  \item Designed a small systems language exposing C's machine model with simplified declaration syntax; built an LLVM-based compiler with interchangeable backends targeting LLVM IR and custom ULM assembly.
  \item Implemented native C interoperability with x86-64 System V ABI lowering, enabling direct use of existing C libraries without wrappers; demonstrated with raylib applications.
  \item Developed \link{https://github.com/michael-lehn/not-abc}{not-abc}, a minimal self-hosting compiler originally implemented in ABC, demonstrating the language's suitability for systems programming and compiler construction.
\end{cvitems}

\cventry{{\mdseries\link{https://github.com/michael-lehn/FLENS}{FLENS}} — Flexible Library for Efficient Numerical Solutions}{C++}
\begin{cvitems}
  \item Built a numerical computing framework combining MATLAB-like abstractions with direct access to optimized BLAS/LAPACK kernels, achieving zero abstraction overhead as verified by benchmarks and generated assembly.
  \item Designed a generic architecture allowing algorithms to operate unchanged on local, distributed, GPU-resident, and implicitly represented matrices.
  \item Used compile-time C++ techniques and expression analysis to map high-level operations directly to optimized numerical kernels and detect performance-critical temporaries.
\end{cvitems}

\cventry{{\mdseries\link{https://github.com/michael-lehn/ulm-on-ice}{ulm-on-iCE}} / {\mdseries\link{https://github.com/michael-lehn/ulm-generator}{ulm-generator}} — FPGA CPU \& ISA Toolchain}{SystemVerilog · C/C++ · Assembly}
\begin{cvitems}
  \item Designed and implemented a custom RISC CPU in SystemVerilog on an iCE40 FPGA, from digital logic and ISA design to assembly programs running on physical hardware.
  \item Defined a declarative ISA description format and built a generator producing an assembler, virtual machine, instruction-level debugger, and documentation from a single specification.
\end{cvitems}

% ===========================================================================
\cvsection{PROFESSIONAL EXPERIENCE}

\cventry{Ulm University — Institute of Numerical Mathematics · Ulm, Germany}{2008 – present}
\cvsubtitle{Tenured Faculty Member (Akademischer Rat auf Lebenszeit), Numerical Mathematics}
\begin{cvitems}
  \item Develop and maintain open-source numerical and systems software in C/C++: high-performance BLAS/LAPACK kernels (ulmBLAS), a numerical computing framework (FLENS), an LLVM-based systems language and compiler (ABC), and ISA/FPGA CPU tooling (ulm-generator / ulm-on-iCE).
  \item Own the full engineering lifecycle for these codebases: architecture, implementation, testing, performance benchmarking and profiling, documentation, and long-term maintenance.
  \item Design and teach courses in high-performance computing, compiler construction, systems programming, numerical mathematics, and scientific computing; the HPC curriculum progresses from memory hierarchy and cache-aware algorithms through SIMD micro-kernels and shared-memory parallelism to MPI and GPU computing.
  \item Built a full-stack computing curriculum spanning logic gates, registers and ALUs, ISA design, assembly, systems programming, and compiler construction, supported by custom-developed software and FPGA hardware.
\end{cvitems}

\cventry{Ulm Center for Scientific Computing (UZWR), Ulm University · Ulm, Germany}{2006 – 2008}
\cvsubtitle{Managing Director}
\begin{cvitems}
  \item Developed complete C++/Qt applications for industrial and research partners.
  \item Translated industrial and interdisciplinary partners' engineering and scientific problems into mathematical models and working numerical software; developed custom solutions where no suitable off-the-shelf software existed.
  \item Among numerous interdisciplinary projects, designed and developed a self-contained Windows application for MRI-based biomechanical analysis, implementing the complete software stack in-house—from FEM and numerical algorithms to data processing and the graphical user interface.
  \item Scoped and evaluated incoming projects, designed computational approaches, and either implemented solutions directly or coordinated them with researchers possessing the required expertise.
\end{cvitems}

\cventry{Ulm University — Institute of Numerical Mathematics · Ulm, Germany}{2002 – 2008}
\cvsubtitle{Research Associate}
\begin{cvitems}
  \item Conducted research and software development in scientific computing and numerical linear algebra, with a focus on high-performance numerical software in C++.
  \item Designed and developed FLENS (Flexible Library for Efficient Numerical Solutions) as the core of my doctoral research on efficient, abstraction-preserving numerical software in C++.
\end{cvitems}

\cventry{University of Southern California — Modeling and Simulation Laboratory · Los Angeles, CA}{1999 – 2001}
\cvsubtitle{Research Assistant / Teaching Assistant}
\begin{cvitems}
  \item Conducted research in applied mathematics and scientific computing with Prof. Chunming Wang in \link{https://dornsife.usc.edu/msl/faculty/}{USC's Modeling and Simulation Laboratory} while completing an M.S. in Applied Mathematics; master's research conducted in collaboration with NASA's Jet Propulsion Laboratory (JPL).
\end{cvitems}

% ===========================================================================
\cvsection{SELECTED INDUSTRY \& RESEARCH COLLABORATIONS}

\cventry{\link{https://www.dlr.de/de/qt}{German Aerospace Center (DLR)} — Real-Time Control / Quantum Technologies}{2025 – present}
\begin{cvitems}
  \item Collaborate with the DLR Institute of Quantum Technologies on control software for a space-ready quantum-memory system based on atomic-gas technology.
  \item Contribute to the architecture and development of the real-time control system and scientist-facing software, using ARTIQ/Sinara with an FPGA-based core device for deterministic control and hardware integration.
  \item Work on real-time state-machine control, experimental parameter management, measurement and imaging acquisition, hardware integration, and FPGA-controlled instrumentation.
\end{cvitems}

\cventry{\link{https://tensor-solutions.com}{Tensor AI Solutions} — HPC / Tensor Networks}{2024}
\begin{cvitems}
  \item Collaborated with Tensor AI Solutions on the design and development of Tartarus, a C++ HPC library for efficient tensor contractions in tree-tensor-network applications for machine learning and QML.
  \item Co-supervised a research-driven student project, contributing substantially to the numerical algorithms, performance-oriented design, and development of the library.
  \item Applied TTGT and Loops-over-GEMM (LoG) contraction schemes, BLAS kernels (GEMM, GEMV, GER, DOT), memory-layout-aware algorithms, and Python-based C++ code generation to optimize tensor contractions on CPUs.
\end{cvitems}

\cventry{\link{https://www.mercedes-benz-techinnovation.com/en}{Daimler TSS} (Daimler Group) — Machine Learning / Recurrent Neural Networks}{2017 – 2018}
\begin{cvitems}
  \item Collaborated closely with Daimler TSS on the mathematical analysis of recurrent neural networks, with a particular focus on training and the vanishing/exploding gradient problem.
  \item Co-supervised a research-driven bachelor's thesis, contributing substantially to the mathematical formulation, analysis, and development of the work.
  \item Applied regression and optimization methods, gradient-based training and backpropagation, and mathematical analysis of feed-forward and Elman recurrent neural networks.
\end{cvitems}

% ===========================================================================
\cvsection{EDUCATION}

\cventry{Ph.D., Numerical Mathematics — Ulm University, Germany}{2008}
\cvsubtitle{Dissertation: FLENS — A Flexible Library for Efficient Numerical Solutions · Advisor: Prof. Karsten Urban}

\cventry{M.S., Mathematics — Ulm University, Germany}{2002}

\cventry{M.S., Applied Mathematics — University of Southern California, Los Angeles, CA}{2001}

\end{document}
'
\documentclass[10pt]{article}

\usepackage[letterpaper,top=32.6bp,bottom=36bp,left=45bp,right=45bp,
            headheight=0bp,headsep=0bp,footskip=0bp]{geometry}
\usepackage{fontspec}
\usepackage[HTML]{xcolor}
\usepackage{enumitem}
\usepackage{lua-ul}
\usepackage{hyperref}

% Pages breaks lines first-fit; TeX's paragraph optimizer would not
\directlua{dofile(kpse.find_file("firstfit.lua", "lua"))}

% Arial if installed, otherwise its metric-compatible clone Arimo.
% No ligatures, and no TeX input mapping (keeps straight quotes as typed).
\IfFontExistsTF{Arial}
  {\setmainfont{Arial}[Ligatures={CommonOff,TeXOff}]}
  {\setmainfont{Arimo}[Ligatures={CommonOff,TeXOff}]}

\definecolor{name}{HTML}{17365D}     % dark blue (name)
\definecolor{tagline}{HTML}{4F81BD}  % light blue (tagline, rule below name)
\definecolor{heading}{HTML}{365F91}  % muted blue (section headings)
\definecolor{rule}{HTML}{808080}     % gray section rules
\definecolor{link}{HTML}{0563C1}     % hyperlinks

\hypersetup{hidelinks, pdftitle={Michael Lehn -- Resume}, pdfauthor={Michael Lehn}}

% Hyperlink in the Pages style: blue and underlined (breakable across lines).
% The underline uses Arial's underline thickness (0.073em regular, 0.105em
% bold), centered 0.1em below the baseline.
\makeatletter
\def\link@thick@m{0.073em}
\def\link@thick@b{0.105em}
\def\link@thick@bx{0.105em}
\newunderlinetype\linkUnderLine[\number\dimexpr1em\relax\f@series]{%
  \leaders\vrule
    height \dimexpr-0.1em+0.5\dimexpr\csname link@thick@\f@series\endcsname\relax\relax
    depth  \dimexpr 0.1em+0.5\dimexpr\csname link@thick@\f@series\endcsname\relax\relax}
\makeatother
\newcommand{\link}[2]{\href{#1}{\textcolor{link}{\linkUnderLine{#2}}}}

% ---------------------------------------------------------------------------
% Font sizes (Pages single line spacing for Arial is 1.145 x font size;
% all body lines use the 10bp line pitch)
\newcommand{\namesize}{\fontsize{22bp}{25.2bp}\selectfont}
\newcommand{\taglinesize}{\fontsize{10.5bp}{16.8bp}\selectfont} % 16.8bp below the name
\newcommand{\headsize}{\fontsize{10.5bp}{11.45bp}\selectfont}
\newcommand{\bodysize}{\fontsize{10bp}{11.45bp}\selectfont}
\newcommand{\smallsize}{\fontsize{9.5bp}{11.45bp}\selectfont}
\newcommand{\contactsize}{\fontsize{9bp}{10.5bp}\selectfont}

% Separator in the header lines
\newcommand{\sep}{\ \ ·\ \ }

% ---------------------------------------------------------------------------
% Vertical spacing: every line uses the 11.45bp line pitch; the extra space
% between blocks is added with \addvspace, so adjacent amounts do not stack.
% The space above an entry depends on what precedes it (\cv@entryskip).
\makeatletter
\newcommand{\cv@entryskip}{4.6bp}

% Section heading: bold blue caps with a 0.5bp gray rule 2.9bp below the
% baseline
\newcommand{\cvsection}[1]{%
  \par\addvspace{6.85bp}%
  {\headsize\bfseries\color{heading}%
    \leavevmode\rlap{\color{rule}\rule[-3.15bp]{\linewidth}{0.5bp}}#1\par}%
  \nopagebreak\addvspace{2.2bp}%
  \gdef\cv@entryskip{4.8bp}%
}

% Entry: bold title, right-aligned small info (technologies or dates)
\newcommand{\cventry}[2]{%
  \par\addvspace{\cv@entryskip}%
  {\bfseries #1}\hfill{\smallsize #2}\kern0.85bp\par % Pages' right tab stop
  \nopagebreak
  \gdef\cv@entryskip{4.1bp}% next entry directly follows (education)
}

% Italic line below an entry (role, dissertation, ...)
\newcommand{\cvsubtitle}[1]{%
  \vspace{0.7bp}%
  {\smallsize\itshape #1\par}%
  \nopagebreak\addvspace{0.9bp}%
}

% Skills paragraph with a bold run-in label
\newcommand{\skill}[2]{%
  \textbf{#1} #2\par\addvspace{2bp}%
}

% Bullet list: bullet at 5bp, text at 15bp from the margin
\newenvironment{cvitems}
  {\begin{itemize}[label=•, leftmargin=15bp, labelwidth=10bp, labelsep=0bp,
                   align=left, topsep=0.55bp, partopsep=0bp, parsep=0bp,
                   itemsep=1.5bp, beginpenalty=10000]}
  {\end{itemize}\gdef\cv@entryskip{4.6bp}}
\makeatother

% ---------------------------------------------------------------------------
\setlength{\parindent}{0bp}
\setlength{\parskip}{0bp}
\setlength{\topskip}{10bp}
\pagestyle{empty}

\begin{document}
\bodysize
\hyphenpenalty=10000 % Pages does not hyphenate this document
\exhyphenpenalty=0
\tolerance=9999      % Pages stretches interword spaces freely in justified text
\emergencystretch=0bp
\frenchspacing       % uniform interword spaces after punctuation, as in Pages
\clubpenalty=10000   % no single lines of a paragraph or bullet at page breaks
\widowpenalty=10000
\makeatletter
\@clubpenalty=\clubpenalty % list items reset \clubpenalty to this value
\makeatother

% ===========================================================================
% Header
{\namesize\addfontfeatures{LetterSpace=1.2}\bfseries\color{name}%
  \rule{0bp}{22bp}% places the baseline 22bp below the top margin
  \rlap{\smash{\color{tagline}\rule[-6.5bp]{\linewidth}{1bp}}}%
  MICHAEL LEHN, Ph.D.\par}
{\centering
  {\taglinesize\addfontfeatures{LetterSpace=7.6}\color{tagline}%
    \spaceskip=\dimexpr\fontdimen2\font+0.8bp\relax % tracking applies to spaces, too
    Scientific Software Engineer\sep C/C++\sep Compilers \& LLVM\sep High-Performance Computing\par}
  \vspace{1.9bp}
  {\contactsize Ulm, Germany\sep Open to relocation to the United States\par}
  {\contactsize
    \link{mailto:michael.lehn@mac.com}{michael.lehn@mac.com}%
    \ifdefined\PublicVersion
      \ \ ·\ \link{https://www.linkedin.com/in/michael-l-2b4a4829b/}{LinkedIn}%
    \else
      \sep +1 (213) 232-7502\sep\link{https://www.linkedin.com/in/michael-l-2b4a4829b/}{LinkedIn}%
    \fi
    \sep\link{https://github.com/michael-lehn}{GitHub}%
    \ ·\ \ \link{https://www.uni-ulm.de/mawi/institut-fuer-numerische-mathematik/institut/mitarbeiter/mlehn/}{Ulm University}\par}
}
\vspace{-0.7bp}

% ===========================================================================
\cvsection{SUMMARY}

Software engineer and mathematician (Ph.D.) with 20+ years of hands-on experience developing performance-critical C/C++ software. Expertise spans numerical algorithms and high-performance linear algebra, memory hierarchy and cache-aware algorithms, SIMD, multithreaded and distributed computing, CPU architecture, compiler construction, and low-level systems programming. Recent projects include an LLVM-based systems programming language with native C interoperability, a minimal self-hosting compiler, and custom CPU architectures implemented on FPGA.\par

% ===========================================================================
\cvsection{TECHNICAL SKILLS}

\skill{Languages:}{C++ (modern, templates/compile-time techniques), C, Python, Assembly (x86-64, AArch64, custom RISC), LLVM IR, SystemVerilog, Java, Perl}
\skill{Performance \& HPC:}{Numerical linear algebra, BLAS, LAPACK, ScaLAPACK, SIMD micro-kernels (SSE3), cache-aware/blocked algorithms, std::thread, OpenMP, MPI, CUDA, GPU and distributed-memory computing, benchmarking and profiling}
\skill{Compilers \& Systems:}{LLVM, compiler construction, code generation, x86-64 System V ABI, C interoperability, ISA design, CPU architecture, FPGA (iCE40, ARTIQ/Sinara), bare-metal programming}
\skill{Tools:}{Linux/Unix, macOS, Git, Qt, CMake, GNU Make, GDB, Valgrind, perf, Intel VTune, NVIDIA Nsight}

% ===========================================================================
\cvsection{SELECTED PROJECTS (OPEN SOURCE)}

\cventry{{\mdseries\link{https://github.com/michael-lehn/ulmBLAS}{ulmBLAS}} — High-Performance BLAS Library}{C++ · Assembly}
\begin{cvitems}
  \item Designed and implemented a BLAS library based on the layered architecture pioneered by BLIS, combining portable blocking/packing algorithms in C++ with small, replaceable architecture-specific micro-kernels.
  \item Reached \textasciitilde97\% of Intel MKL DGEMM performance with BLIS micro-kernels; wrote custom SSE3 micro-kernels and a plug-in interface for externally optimized kernels.
  \item Drove blocked LAPACK routines (e.g., DGETRF) at performance comparable to vendor LAPACK implementations.
  \item Authored the \link{https://github.com/michael-lehn/gemm-tutorial}{GEMM Tutorial} on optimizing matrix multiplication from pure C to SSE micro-kernels; referenced by the \link{https://github.com/flame/how-to-optimize-gemm}{FLAME project's GEMM optimization resources}.
  \item Used in \link{https://www.meta.com/de/en/legal/quest/third-party-notices/14/}{Meta Quest software}; listed in Meta’s official third-party software notices.
\end{cvitems}

\cventry{{\mdseries\link{https://github.com/michael-lehn/abc-llvm}{ABC}} ("A Better C") — Systems Language \& LLVM Compiler}{C++ · LLVM}
\begin{cvitems}
  \item Designed a small systems language exposing C's machine model with simplified declaration syntax; built an LLVM-based compiler with interchangeable backends targeting LLVM IR and custom ULM assembly.
  \item Implemented native C interoperability with x86-64 System V ABI lowering, enabling direct use of existing C libraries without wrappers; demonstrated with raylib applications.
  \item Developed \link{https://github.com/michael-lehn/not-abc}{not-abc}, a minimal self-hosting compiler originally implemented in ABC, demonstrating the language's suitability for systems programming and compiler construction.
\end{cvitems}

\cventry{{\mdseries\link{https://github.com/michael-lehn/FLENS}{FLENS}} — Flexible Library for Efficient Numerical Solutions}{C++}
\begin{cvitems}
  \item Built a numerical computing framework combining MATLAB-like abstractions with direct access to optimized BLAS/LAPACK kernels, achieving zero abstraction overhead as verified by benchmarks and generated assembly.
  \item Designed a generic architecture allowing algorithms to operate unchanged on local, distributed, GPU-resident, and implicitly represented matrices.
  \item Used compile-time C++ techniques and expression analysis to map high-level operations directly to optimized numerical kernels and detect performance-critical temporaries.
\end{cvitems}

\cventry{{\mdseries\link{https://github.com/michael-lehn/ulm-on-ice}{ulm-on-iCE}} / {\mdseries\link{https://github.com/michael-lehn/ulm-generator}{ulm-generator}} — FPGA CPU \& ISA Toolchain}{SystemVerilog · C/C++ · Assembly}
\begin{cvitems}
  \item Designed and implemented a custom RISC CPU in SystemVerilog on an iCE40 FPGA, from digital logic and ISA design to assembly programs running on physical hardware.
  \item Defined a declarative ISA description format and built a generator producing an assembler, virtual machine, instruction-level debugger, and documentation from a single specification.
\end{cvitems}

% ===========================================================================
\cvsection{PROFESSIONAL EXPERIENCE}

\cventry{Ulm University — Institute of Numerical Mathematics · Ulm, Germany}{2008 – present}
\cvsubtitle{Tenured Faculty Member (Akademischer Rat auf Lebenszeit), Numerical Mathematics}
\begin{cvitems}
  \item Develop and maintain open-source numerical and systems software in C/C++: high-performance BLAS/LAPACK kernels (ulmBLAS), a numerical computing framework (FLENS), an LLVM-based systems language and compiler (ABC), and ISA/FPGA CPU tooling (ulm-generator / ulm-on-iCE).
  \item Own the full engineering lifecycle for these codebases: architecture, implementation, testing, performance benchmarking and profiling, documentation, and long-term maintenance.
  \item Design and teach courses in high-performance computing, compiler construction, systems programming, numerical mathematics, and scientific computing; the HPC curriculum progresses from memory hierarchy and cache-aware algorithms through SIMD micro-kernels and shared-memory parallelism to MPI and GPU computing.
  \item Built a full-stack computing curriculum spanning logic gates, registers and ALUs, ISA design, assembly, systems programming, and compiler construction, supported by custom-developed software and FPGA hardware.
\end{cvitems}

\cventry{Ulm Center for Scientific Computing (UZWR), Ulm University · Ulm, Germany}{2006 – 2008}
\cvsubtitle{Managing Director}
\begin{cvitems}
  \item Developed complete C++/Qt applications for industrial and research partners.
  \item Translated industrial and interdisciplinary partners' engineering and scientific problems into mathematical models and working numerical software; developed custom solutions where no suitable off-the-shelf software existed.
  \item Among numerous interdisciplinary projects, designed and developed a self-contained Windows application for MRI-based biomechanical analysis, implementing the complete software stack in-house—from FEM and numerical algorithms to data processing and the graphical user interface.
  \item Scoped and evaluated incoming projects, designed computational approaches, and either implemented solutions directly or coordinated them with researchers possessing the required expertise.
\end{cvitems}

\cventry{Ulm University — Institute of Numerical Mathematics · Ulm, Germany}{2002 – 2008}
\cvsubtitle{Research Associate}
\begin{cvitems}
  \item Conducted research and software development in scientific computing and numerical linear algebra, with a focus on high-performance numerical software in C++.
  \item Designed and developed FLENS (Flexible Library for Efficient Numerical Solutions) as the core of my doctoral research on efficient, abstraction-preserving numerical software in C++.
\end{cvitems}

\cventry{University of Southern California — Modeling and Simulation Laboratory · Los Angeles, CA}{1999 – 2001}
\cvsubtitle{Research Assistant / Teaching Assistant}
\begin{cvitems}
  \item Conducted research in applied mathematics and scientific computing with Prof. Chunming Wang in \link{https://dornsife.usc.edu/msl/faculty/}{USC's Modeling and Simulation Laboratory} while completing an M.S. in Applied Mathematics; master's research conducted in collaboration with NASA's Jet Propulsion Laboratory (JPL).
\end{cvitems}

% ===========================================================================
\cvsection{SELECTED INDUSTRY \& RESEARCH COLLABORATIONS}

\cventry{\link{https://www.dlr.de/de/qt}{German Aerospace Center (DLR)} — Real-Time Control / Quantum Technologies}{2025 – present}
\begin{cvitems}
  \item Collaborate with the DLR Institute of Quantum Technologies on control software for a space-ready quantum-memory system based on atomic-gas technology.
  \item Contribute to the architecture and development of the real-time control system and scientist-facing software, using ARTIQ/Sinara with an FPGA-based core device for deterministic control and hardware integration.
  \item Work on real-time state-machine control, experimental parameter management, measurement and imaging acquisition, hardware integration, and FPGA-controlled instrumentation.
\end{cvitems}

\cventry{\link{https://tensor-solutions.com}{Tensor AI Solutions} — HPC / Tensor Networks}{2024}
\begin{cvitems}
  \item Collaborated with Tensor AI Solutions on the design and development of Tartarus, a C++ HPC library for efficient tensor contractions in tree-tensor-network applications for machine learning and QML.
  \item Co-supervised a research-driven student project, contributing substantially to the numerical algorithms, performance-oriented design, and development of the library.
  \item Applied TTGT and Loops-over-GEMM (LoG) contraction schemes, BLAS kernels (GEMM, GEMV, GER, DOT), memory-layout-aware algorithms, and Python-based C++ code generation to optimize tensor contractions on CPUs.
\end{cvitems}

\cventry{\link{https://www.mercedes-benz-techinnovation.com/en}{Daimler TSS} (Daimler Group) — Machine Learning / Recurrent Neural Networks}{2017 – 2018}
\begin{cvitems}
  \item Collaborated closely with Daimler TSS on the mathematical analysis of recurrent neural networks, with a particular focus on training and the vanishing/exploding gradient problem.
  \item Co-supervised a research-driven bachelor's thesis, contributing substantially to the mathematical formulation, analysis, and development of the work.
  \item Applied regression and optimization methods, gradient-based training and backpropagation, and mathematical analysis of feed-forward and Elman recurrent neural networks.
\end{cvitems}

% ===========================================================================
\cvsection{EDUCATION}

\cventry{Ph.D., Numerical Mathematics — Ulm University, Germany}{2008}
\cvsubtitle{Dissertation: FLENS — A Flexible Library for Efficient Numerical Solutions · Advisor: Prof. Karsten Urban}

\cventry{M.S., Mathematics — Ulm University, Germany}{2002}

\cventry{M.S., Applied Mathematics — University of Southern California, Los Angeles, CA}{2001}

\end{document}
